\chapter*{导\quad 论}
\phantomsection
\addcontentsline{toc}{chapter}{导\quad 论}
\markboth{导\quad 论}{导\quad 论}

\noindent
一九六三年春，莫斯科，苏联科学院远东委员会设于一栋旧楼三层的档案室，一名读汉文的研究生被叫去看三只木箱。编目卡上的类别写作「远东出土文书」，箱内是三十七片长约二寸至九寸、宽约一指的薄片，多为金属，间有陶瓷，以棉纸逐片隔开。片上有刻字，字为小篆。

此后的六十年里，这批材料经过三次编目、一次转移、两次被列入清退名单，始终没有人调阅。使它搁置的原因与保密无关，只是没有人知道应当把它归入哪一类：从材质看是器物，从字迹看是文书，从出土地点看则不属于任何一个国家的考古区划。一九九二年基金会接管远东委员会遗留卷宗之后，它被编入 \textsc{фв} 系列，编号自 \textsc{фв}-1963-11 起，与一批来历相近、彼此并无关联的散件混在一处。

本书所处理的，就是这三十七片。

\section*{一、材料的性质}

关于这批材料的出处，档案只有一行记载：

\begin{center}
\small
\textsc{фв}-1963-11\quad 月球正面，\textsc{mare-tranq-07} 基地，第 4 号居住环，西侧储物间
\end{center}

\textsc{mare-tranq-07} 即《战争的雷霆》第一章所考的静海七号，也是汉译「广寒宫」所指的那一处。同批材料内另存两张照片，摄于取走残片的那个房间，画面所示西壁架数与简七的自记相合，可证三十七片确出自该处。

这批材料在史料学上的位置需要说明。它属于原件，并非从他书引文中钩沉出来的佚文：字口清晰，可辨部分约占十分之六七，残损处蚀去的是字，不是载体。它同样不属于出土文书，因为出土地点不在中国境内，而在月球正面的月海之中。这两点使它与传统的辑佚工作区别开来。

苏联方面在一九六三年曾对同类材料作过一次成分分析，报告编号 \textsc{фв}-1963-07，结论部分只留下一句「来源不明，暂不归类」。这句话写在编目卡上，不在原件上，但它准确地概括了此后六十年间的处境。

\section*{二、材料的内容}

三十七片构不成一部完整的书。它们是一处设施在长时段内陆续刻写的簿籍与规程：居住环的形制、气路与水路的分布、人口的编制、口粮的定额、名册的更替、事故的数目，以及撤离前的清点与退役。

其中分量最重的是卷三的十二片。它们记的是公元前二百二十一年前后的一次清点，与《战争的雷霆》第一章所引「第 2187 号卷宗第 121 叶，设施退役清单，共十九项」同出一源。那一处只引了第七项，本书给出全文。

就内容而言，这批材料所记的是一处基地的日常行政，而非任何宏大事件。三十七片记下的是岁一检的管道、岁一焚的名册、日三合的定额，以及最后一年里逐条执行的十九项处置。人名、战争、决定，均无记载。唯一一处带有动作的记载，是简三十一「闭者七人，回首一次」。

就材料所及，可以为此城勾出一个轮廓。它在月球正面静海西缘，直径约一点二公里，六环相套，居民十一万九千人。居民不是地球人，其体质与地上人不属于同一类群：不老化，不生育，除外伤与极大量的污秽之外不会死亡。城的洁净不靠天然，靠一套藏在外墙之内的系统常年运转，而系统的能源取自地下。公元前二百二十一年，地下热源降至临界以下，全城分六批离城，城随即封存。

以上数语皆为下编各章的结论，此处提前写出，只为使读者在进入释文之前知道这批字记的是什么。各条的依据与限定见下编，尤其是第二章、第六章与第七章。凡传世文献与此城相关之事，本书一概不涉，其故见下编前言。

这一特征既是材料的局限，也是它的用处。行政记录不修辞，不表彰，不隐恶；它记下的是当时的人认为需要留下来的东西。至于什么被认作需要留下来，本身就是一项可以研究的对象。

\section*{三、体例}

本书分上下两编。

上编为释文，按内容分四卷：营建、岁计、撤离、后题，另附存疑八条。每片三段：释文照原件行款，出处记编号、尺寸与材质，校记则处理字词。校记不作论断，凡涉及释读之判断，一律移入下编。

下编为考释，共八章。第一章论月之民的体质与制度，第二章据简六所记尺度复原城的形制，第三章论净秽制度，第四章只处理两个字，即简二十二第十三项「药，弃」，第五章论净秽山与月兔，第六章处理岁计与人口的数目，第七章还原撤离的经过，第八章论名与籍。

两编的读法不同。上编可以跳读，按下编索引所举的简次查阅即可；下编要从头读起，因为结论是逐层建立起来的。

凡例另列于上编之前，不再重复。有一件事要在这里交代：原件所刻为小篆，释文一律改作今字，这是整理工作的通例，不是原件的字形。书中的插图同样照释文摹写。要见原字形，只能看影印本；影印本拟随第三版附出，此次未及。

\section*{四、方法与限度}

本书的编次出自笔者。原件的编号与摆放的次序均已无从确认，照片上能辨的只剩层数。现在的编次依据材质、字口与内容三类线索的交叉，其中任何一类单独成立都不足以定序。换一个人来做，次序会不同。

释文的情况与此相同。可辨的十之六七中，约有一成为笔者所补；凡补字一概用方框标出，校记中逐条说明依据。读者可以不同意，也可以另补，原件尚在，谁都可以重来。

数字方面，凡能互相校验的都已校过。校验不上的写在正文里，不加隐去，其中最显著的是八千之数：完工时十五万人（见《战争的雷霆》第一章第三节），六批一共下去十二万人，两万年间少了三万；简十九记事故二万二千，两数相减余八千。这八千人在残简中无一字记载。本书第六章列出三种解释，倾向第一种，并声明没有依据。

最后是年表的问题。城中自有纪年，与中原不同，无法换算。残简中提及城中纪年者三处（十四年、十九年、二十三年），另有一处作「元年」。前者是否与后者同属一纪年，无从判定，故附录二将其分列。

\section*{五、自白}

笔者不认为本书的编次是正确的，只认为它是一种必要的编次。

三十七片放在一间屋子里六十年，没有人调阅。整理成书并不指望改变这一点。这批材料大概率仍会留在那间屋子里，至多换一个编号。若要说有什么用处，那就是把这三十七片刻的字，从金属与陶瓷上移到纸上。纸不褪色，也不腐烂，只要不被烧掉，它就一直待在那里，等人来翻。

\begin{flushright}
{\small 巴兰尼科夫}\quad{\footnotesize\color{pencil} 二〇二四年六月，于莫斯科}
\end{flushright}
